\chapter{Las reglas de Feymnan }\label{cap.2}
\markboth{Teoría débil y modelo estándar}{Teoría débil y modelo estándar}
% Colocar el resumen del capítulo

\section{Las reglas de Feymnan para una teoría de juguete} \renewcommand{\thefootnote}{\arabic{footnote}} % numeración por símbolos
\lhead[\thepage]{\thesection. Las reglas de Feymnan para una teoría de juguete}

En la Sección 6.2 aprendimos cómo calcular tasas de desintegración y secciones transversales de dispersión, en términos de la amplitud \( \mathcal{M} \) para el proceso en cuestión. Ahora te mostraré cómo determinar \( \mathcal{M} \) misma, usando las "reglas de Feynman" para evaluar los diagramas relevantes. Podríamos ir directamente a un sistema de la "vida real", como la electrodinámica cuántica, con electrones y fotones interactuando a través del vértice primitivo:

\tikzfeynmanset{compat=1.1.0}
\begin{figure}[h!]
    \centering
    \begin{tikzpicture}
      \begin{feynman}
        % Vértices
        \vertex (a) {\(e\)};
        \vertex [below=2cm of a] (b) {\(e\)};
        \vertex [right=2cm of a] (f1);
        \vertex [right=2cm of b] (f2);
        \vertex [right=1.5cm of $(a)!0.5!(b)$] (c);

        % Líneas
        \diagram* {
          (a) -- [fermion] (c) -- [fermion] (b),
          (c) -- [photon] (f1) [particle=\(\gamma\)],
        };
      \end{feynman}
    \end{tikzpicture}
    \caption{Diagrama de Feynman para la emisión de un fotón por un electrón.}
    \label{fig:emision_foton}
\end{figure}


Esta es la aplicación original, la más importante y la mejor entendida de la técnica de Feynman. 
Desafortunadamente, involucra complicaciones adicionales (debido al hecho de que el electrón tiene espín 
$\tfrac{1}{2}$ y el fotón espín $1$) que no tienen nada que ver con el cálculo de Feynman en sí. 
En el Capítulo 7 mostraré cómo manejar partículas con espín, pero por el momento no quiero confundir 
el tema, así que voy a introducir una ``teoría de juguete'', que no pretende representar el mundo real, 
sino que servirá para ilustrar el \textit{método}, con un mínimo de equipaje innecesario.

\bigskip

Imaginemos un mundo en el que solo existen tres tipos de partículas —llamémoslas $A, B, C$— con masas 
$m_A, m_B, m_C$. Todas tienen espín $0$ y cada una es su propia antipartícula. Existe un único vértice 
primitivo mediante el cual las tres partículas interactúan:

\begin{figure}[h!]
    \centering
    \begin{tikzpicture}
      \begin{feynman}
        % Vértices
        \vertex (a) at (210:2) {\(A\)};
        \vertex (b) at (0:2) {\(B\)};
        \vertex (c) at (120:2) {\(C\)};
        \vertex (v) at (0,0);

        % Líneas
        \diagram* {
          (a) -- (v) -- (b),
          (c) -- (v),
        };
      \end{feynman}
    \end{tikzpicture}
    \caption{Vértice primitivo de interacción entre las partículas $A$, $B$ y $C$.}
\end{figure}

Supongamos que $A$ es la más pesada de las tres, y de hecho tiene más masa que 
$B$ y $C$ combinadas, de modo que puede decaer en $B + C$. 
El diagrama de más bajo orden que describe esta desintegración es:

\begin{figure}[h!]
    \centering
    \begin{tikzpicture}
      \begin{feynman}
        % Vértices
        \vertex (a) at (270:2) {\(A\)};
        \vertex (b) at (150:2) {\(B\)};
        \vertex (c) at (30:2) {\(C\)};
        \vertex (v) at (0,0);

        % Líneas
        \diagram* {
          (a) -- (v) -- (b),
          (v) -- (c),
        };
      \end{feynman}
    \end{tikzpicture}
    \caption{Desintegración de la partícula $A$ en $B$ y $C$ al orden más bajo.}
\end{figure}


